% GENERATED FILE. Edit canonical lessons, then run npm run generate:books.
% canonical-source-hash: fnv1a64:da5af949620c8c8c
% canonical-lessons: TE-C67-wall, TE-C67-roof, TE-C67-threshold, TE-C67-step, TE-C67-window

\chapter{The House Itself}
\label{ch:te-the-house-itself}
\hlchaptermodality{67}

\begin{chapteropening}
\textbf{By the end of this chapter:} \emph{I can name the parts a Telugu house is built from, and greet the people inside it.}

Complete TE-C67-window and demonstrate the chapter promise.
\end{chapteropening}

\section[gōḍa]{\te{గోడ} (gōḍa) — a wall}

\label{lesson:TE-C67-wall}



\begin{quote}
Before the new one: what did \emph{paṇṭa} mean?
\end{quote}

\subsection*{You'll want to know: \te{గోడ}}
\textbf{\te{గోడ}} (\emph{gōḍa}) — \textquotedblleft{}a wall\textquotedblright{}.

\te{గోడ} is a wall: the built side of a room. A village house puts up four of them in mud or in brick, and everything after that — the \te{తలుపు}, the \te{కుర్చీ} pushed back against one of them — is placed by where they stand.

The word is inherited, and it does the same figurative work English asks of a wall. Somebody who will not listen is a \te{గోడ}, and a family with nothing hidden has no \te{గోడ} inside it.

Whitewashing the \te{గోడ} before a festival is its own small season in a Telugu house. A guest who arrives to a fresh one has been told they were expected, and the right thing to say at that moment is \te{నమస్కారం}.

\begin{grammarlens}[title={What you've built}]
One. What a house is measured by before anything is put inside it.
\end{grammarlens}

\subsection*{Your turn}
Say these aloud:

\begin{itemize}
\raggedright
  \item \emph{gōḍa}
  \item it once more, slowly
  \item \emph{gōḍa}, then \emph{namaskāram}, as though greeting somebody at a freshly painted one
\end{itemize}

\begin{tcolorbox}[breakable,colback=teal!4,colframe=teal!35!black,title={Before you move on}]
What does \te{గోడ} mean? (\textquotedblleft{}a wall\textquotedblright{}.) Say it once more.
\end{tcolorbox}
\section[kappu]{\te{కప్పు} (kappu) — a roof}

\label{lesson:TE-C67-roof}



\begin{quote}
Before the new one: what did \emph{gōḍa} mean?
\end{quote}

\subsection*{You'll want to know: \te{కప్పు}}
\textbf{\te{కప్పు}} (\emph{kappu}) — \textquotedblleft{}a roof\textquotedblright{}.

\te{కప్పు} is the roof — the thing that covers. It comes from the same stem as the everyday word for covering something over, so a lid on a \te{గిన్నె} and the roof over a room are one idea in Telugu.

Village roofs were \te{గడ్డి} over beams for centuries and are tile or concrete now. The material changed and the word did not, which is usually a sign of a very old one.

To be under somebody's \te{కప్పు} is to be their guest, and it binds the person underneath as much as the person above. \te{నమస్కారం} comes first, and everything else follows it.

\begin{grammarlens}[title={What you've built}]
Two: \te{గోడ} and \te{కప్పు}. Sides, and a lid.
\end{grammarlens}

\subsection*{Your turn}
Say these aloud:

\begin{itemize}
\raggedright
  \item \emph{kappu}
  \item it once more, slowly
  \item \emph{gōḍa}, then \emph{kappu}, and say which one holds up the other
\end{itemize}

\begin{tcolorbox}[breakable,colback=teal!4,colframe=teal!35!black,title={Before you move on}]
What does \te{కప్పు} mean? (\textquotedblleft{}a roof\textquotedblright{}.) Say it once more.
\end{tcolorbox}
\section[gaḍapa]{\te{గడప} (gaḍapa) — a threshold}

\label{lesson:TE-C67-threshold}



\begin{quote}
Before the new one: what did \emph{kappu} mean?
\end{quote}

\subsection*{You'll want to know: \te{గడప}}
\textbf{\te{గడప}} (\emph{gaḍapa}) — \textquotedblleft{}a threshold\textquotedblright{}.

\te{గడప} is the threshold: the raised strip of stone or wood under the \te{తలుపు} that you step over on the way in.

It is the most watched piece of a Telugu house. The \te{ముగ్గు} is drawn just outside it, water is thrown across it at \te{ఉదయం}, and nobody steps on it if they can help it.

Because it marks in from out, \te{గడప} also does the work of the word household. Asking whose \te{గడప} somebody comes from is asking about their family, not about their doorstep.

\begin{grammarlens}[title={What you've built}]
Three. The line you cross to be inside.
\end{grammarlens}

\subsection*{Your turn}
Say these aloud:

\begin{itemize}
\raggedright
  \item \emph{gaḍapa}
  \item it once more, slowly
  \item \emph{gaḍapa}, then \emph{muggu}, and say which one is drawn beside the other
\end{itemize}

\begin{tcolorbox}[breakable,colback=teal!4,colframe=teal!35!black,title={Before you move on}]
What does \te{గడప} mean? (\textquotedblleft{}a threshold\textquotedblright{}.) Say it once more.
\end{tcolorbox}
\section[meṭṭu]{\te{మెట్టు} (meṭṭu) — a step}

\label{lesson:TE-C67-step}



\begin{quote}
Before the new one: what did \emph{gaḍapa} mean?
\end{quote}

\subsection*{You'll want to know: \te{మెట్టు}}
\textbf{\te{మెట్టు}} (\emph{meṭṭu}) — \textquotedblleft{}a step\textquotedblright{}.

\te{మెట్టు} is a step — one of a flight, cut in stone or built up in mud.

Telugu counts them and remembers them. The steps down to a \te{చెరువు} are named in some villages, and a house with a \te{మెట్టు} at its \te{గడప} is a house standing a little above the street.

The word carries rank as well. A person who has risen has climbed a \te{మెట్టు}, which is how a stair turns into a career with nobody explaining it.

\begin{grammarlens}[title={What you've built}]
Four. Something to climb before you cross.
\end{grammarlens}

\subsection*{Your turn}
Say these aloud:

\begin{itemize}
\raggedright
  \item \emph{meṭṭu}
  \item it once more, slowly
  \item \emph{meṭṭu}, then \emph{gaḍapa}, in the order your feet meet them
\end{itemize}

\begin{tcolorbox}[breakable,colback=teal!4,colframe=teal!35!black,title={Before you move on}]
What does \te{మెట్టు} mean? (\textquotedblleft{}a step\textquotedblright{}.) Say it once more.
\end{tcolorbox}
\section[kiṭikī]{\te{కిటికీ} (kiṭikī) — a window}

\label{lesson:TE-C67-window}



\begin{quote}
Before the new one: what did \emph{meṭṭu} mean?
\end{quote}

\subsection*{You'll want to know: \te{కిటికీ}}
\textbf{\te{కిటికీ}} (\emph{kiṭikī}) — \textquotedblleft{}a window\textquotedblright{}.

\te{కిటికీ} is a window. It is a borrowing, ultimately from Portuguese, and it arrived carrying the shape of window it names: a framed opening in a \te{గోడ}, rather than the small unglazed gap older houses left for \te{గాలి}.

Telugu kept the borrowed word and the older habit side by side. A \te{కిటికీ} is opened at \te{ఉదయం}, shut against the sun in the middle of the day, and opened again by \te{సాయంత్రం}.

A house with a \te{గోడ}, a \te{కప్పు}, a \te{గడప}, a \te{మెట్టు} and a \te{కిటికీ} is a whole house. Whoever is standing at any one of them is somebody you can greet.

\begin{grammarlens}[title={What you've built}]
Five: \te{గోడ}, \te{కప్పు}, \te{గడప}, \te{మెట్టు}, \te{కిటికీ}. A house you could walk up to and greet.
\end{grammarlens}

\subsection*{Your turn}
Say these aloud:

\begin{itemize}
\raggedright
  \item \emph{kiṭikī}
  \item it once more, slowly
  \item all five in order, then \emph{namaskāram} at the last one
\end{itemize}

\begin{tcolorbox}[breakable,colback=teal!4,colframe=teal!35!black,title={Before you move on}]
What does \te{కిటికీ} mean? (\textquotedblleft{}a window\textquotedblright{}.) Say it once more.
\end{tcolorbox}
